% GDZ PPN236018647,native517–519upper;HKUST518–520.
\clearpage\SourceLeaf{517}\NoSourceHead
\pdfbookmark[1]{Abdruck des Tagebuchs mit Erläuterungen: 1–6}{diary-printed-1}
\begin{center}\bfseries
ABDRUCK DES TAGEBUCHS (NOTIZENJOURNALS)\\
MIT ERLÄUTERUNGEN.
\end{center}
\begingroup
\newcommand{\DiaryCeres}{\tikz[baseline=-.1ex,x=1.8ex,y=1.8ex,line width=.06em]{\draw (.05,.65) .. controls (.28,1.13) and (.85,1.03) .. (.8,.55) .. controls (.8,.12) and (.15,.12) .. (.03,.42);\draw (.1,.4)--(.1,-.15);\draw (-.08,.05)--(.34,.05);}}
\newcommand{\DiaryPallas}{\tikz[baseline=-.1ex,x=1.8ex,y=1.8ex,line width=.06em]{\draw (0,.45)--(.35,1)--(.7,.45)--cycle;\draw (.35,.45)--(.35,-.15);\draw (.12,.1)--(.58,.1);}}
\NumberedItem{1}
\indent Principia quibus innititur sectio circuli, ac divisibilitas eiusdem geometrica in septendecim partes etc.
\par\noindent\hfill[1796] Mart. 30.\quad Brunsv[igae]

\medskip\begingroup\small
\indent Das gleiche Datum Werke I, S. 476 (handschriftliche Bemerkung zum art. 365. der \textit{Disquis. Arithm.}). Vergl. die oben S. 3 abgedruckte Anzeige im Intelligenzblatt der allgemeinen Litteraturzeitung Nr. 66 vom 1. Juni 1796. Zur Geschichte der Entdeckung vergl. den oben S. 121 abgedruckten Brief an \textsc{Gerling} vom 6. Januar 1819, ferner zwei Stellen aus Briefen von W. \textsc{Bolyai}, die nach dem Abdruck in dem \textit{Briefwechsel zwischen C. F. Gauss und W. Bolyai}, herausgegeben von P. \textsc{Stäckel} und F. \textsc{Schmidt}, Leipzig 1899, S. 144 und 152 hier wiedergegeben seien. Am 12. April 1855 schreibt \textsc{Bolyai} an den Wiener Astronomen \textsc{Kreil} »und (im \textit{Disquis. Arithm.} pag. 662) \(\cos\frac{P}{17}=\ldotp\,\ldotp\,\ldotp\) war noch vor 1797 erfunden, weil er es bei mir dann an einem Abende in Göttingen sagte« und am 13. Juli 1855 an W. \textsc{Sartorius v. Waltershausen}: »\(\ldotp\,\ldotp\,\ldotp\) nur einmal sah ich an ihm eine mäßige Freude, wo er die kleine Tafel, auf welcher er das 17eck \textit{Disquis. Arithm.} p. 662 berechnet hat, zum Andenken mir gab, und [die] ich nun als etwas interessantes [nach Göttingen] hinschicke«.

\indent Auf S. 77 der 1801 begonnenen Scheda Af hat \textsc{Gauss} die folgende Zusammenstellung für ihn wichtiger Daten eingetragen:
\begin{center}
\setlength{\fboxsep}{9pt}\fbox{\begin{tabular}{rlll}
1. & Jan. & 1801 & \DiaryCeres\ entdeckt\\
19. & Febr. & 1803 & \DiaryPallas\ wiedergefunden\\
28. & März & 1802 & \DiaryPallas\ entdeckt\\
30. & März & 1796 & Construction des 17 Ecks\\
30. & April & 1777 & \(\ast\)\\
7. & Dec. & 1801 & \DiaryCeres\ wiedergefunden
\end{tabular}}
\end{center}
\noindent\centerline{Es bedeutet hier \DiaryCeres\ Ceres, \DiaryPallas\ Pallas.}
\par\noindent\hfill\textsc{Klein. Schlesinger.}
\par\endgroup\endgroup

\clearpage\SourceLeaf{518}\ArithmeticHead{489}{R}{1796 MART. 30—1796 APR. 29.}
\NumberedItem{2}
\indent Numerorum primorum non omnes numeros infra ipsos residua quadratica esse posse demonstratione munitum.
\par\noindent\hfill[1796] Apr. 8.\quad Ibid. [Brunsvigae]

\medskip\begingroup\small
\indent Die Zeitangabe stimmt mit einer Werke I, S. 475 wiedergegebenen handschriftlichen Bemerkung von \textsc{Gauss} überein, wo die Worte des art. 130. der \textit{Disquisitiones arithmeticae} (Werke I, S. 98, 99) »Postquam rigorose demonstravimus, quemvis numerum primum formae \(4n+1\), et positive et negative acceptum, alicuius numeri primi ipso minoris non-residuum esse \(\ldotp\,\ldotp\,\ldotp\)« mit folgender Bemerkung begleitet werden: »hanc demonstrationem deteximus 1796 Apr. 8.« Im Anschluß an diesen Satz gibt \textsc{Gauss} im art. 131. der \textit{Disquisitiones arithmeticae} den ersten seiner sechs Beweise des quadratischen Reziprozitätsgesetzes, bezüglich dessen es Werke I, S. 475 weiter heißt: »Theorema fundamentale per inductionem detectum 1795 Martio. Demonstratio prima, quae in hac sectione traditur, inventa 1796 Apr.«. Der Wortlaut der Tagebuchnotiz Nr. 2 stimmt aber nicht mit dem Satze des art. 130. sondern mit dem sehr viel einfacheren des art. 96. der \textit{Disquisitiones arithmeticae} (Werke I, S. 74) überein. \textsc{Gauss} muß also entweder die Beweise \textls[120]{beider} Sätze am 8. April 1796 gefunden haben, oder die Tagebuchnotiz ist verschrieben und sollte lauten: Numeros primos non omnium numerorum infra ipsos etc. Das letztere ist darum anzunehmen, weil der im Texte der Tagebuchnotiz angezeigte Beweis dafür, daß quadratische Reste einer Primzahl nicht alle unter ihr liegenden Zahlen sein können, \textsc{Gauss} sicher schon bekannt war, als er (März 1795) das Fundamentaltheorem durch Induktion fand, während er den Beweis des im art. 130. enthaltenen Satzes, daß wenigstens eine Primzahl \(p'<p=4n+1\) vorhanden sei, von der \(p\) Nichtrest ist, wie Werke II, S. 4 (1808) berichtet wird, ein ganzes Jahr lang (nämlich März 1795 bis April 1796) trotz angestrengtester Arbeit nicht zu zwingen vermochte.
\par\noindent\hfill\textsc{Klein. Bachmann.}
\par\endgroup

\NumberedItem{3}
\indent Formulae pro cosinibus angulorum peripheriae submultiplorum expressionem generaliorem non admittent nisi in duab[us] periodis.
\par\noindent\hfill[1796] Apr. 12.\quad Ibid. [Brunsvigae]

\NumberedItem{4}
\indent Amplificatio normae residuorum ad residua et mensuras non indivisibiles.
\par\noindent\hfill[1796] Apr. 29.\quad Götting[ae]

\medskip\begingroup\small
\indent Diese Notiz bezieht sich auf das verallgemeinerte quadratische Reziprozitätsgesetz, das \textsc{Gauss} später im art. 133. der \textit{Disquisitiones arithmeticae}, Werke I, S. 101, entwickelt hat; wegen des Datums vergl. die Bemerkung zu diesem Artikel, Werke I, S. 476.
\par\noindent\hfill\textsc{Klein. Bachmann.}
\par\endgroup
\par\vfill\noindent\small X 1.\hfill62\hspace*{.43\linewidth}\normalsize

\clearpage\SourceLeaf{519}\ArithmeticHead{490}{L}{TAGEBUCH MIT ERLÄUTERUNGEN.}
\NumberedItem{5}
\indent Numeri cuiusvis divisibilitas varia in binos primos.
\par\noindent\hfill[1796] Mai. 14.\quad Gött[ingae]

\medskip\begingroup\small
\indent Es ist hier offenbar der Satz gemeint, daß jede gerade Zahl als die Summe zweier Primzahlen dargestellt werden kann. Man bezeichnet diesen empirischen Satz gemeinhin als \textsc{Goldbach}schen, unter Berufung auf die Briefe \textsc{Goldbachs} und \textsc{Eulers} vom 7. und 30. Juni 1742\footnote{Siehe P. H. \textsc{Fuss}, \textit{Correspondance mathématique et physique etc.} I, St. Petersburg 1843, S. 127 und 135.}. Anscheinend geht diese Bezeichnung auf eine Notiz in den Nouvelles Annales de Mathématiques \textbf{14}, 1855, S. 117 zurück, die von dem damaligen Herausgeber \textsc{Terquem} herrührt\begingroup\renewcommand{\thefootnote}{**)}\footnote{S. 293 desselben Bandes findet sich eine \textit{Note sur le théorème de Goldbach} von \textsc{Desboves}.}\endgroup. Im 18. Bande (1859) derselben Zeitschrift, S. 2 des Bulletin de Bibliographie etc., bemerkt \textsc{Terquem} jedoch, daß sich der gedachte Satz auf S. 379 der 3. Ausgabe von E. \textsc{Warings} \textit{Meditationes algebraicae}\begingroup\renewcommand{\thefootnote}{***)}\footnote{Editio tertia recensita et aucta, Cantabrigiae 1782. Die erste Ausgabe war 1770 erschienen.}\endgroup findet, wo es heißt: »Omnis par numerus constat e duobus primis numeris et omnis impar numerus vel est primus numerus, vel constat e tribus primis numeris etc.«\begingroup\renewcommand{\thefootnote}{\ensuremath{\dagger})}\footnote{Vergl. auch die Aufsätze von \textsc{Catalan} und von \textsc{Eneström}, Bulletino di bibliografia e di storia delle scienze matem. e fisiche, pubbl. da B. \textsc{Boncompagni} \textbf{17}, 1884, S. 457, 458.}\endgroup.
\par\noindent\hfill\textsc{Stäckel. Schlesinger.}
\par\endgroup

\NumberedItem{6}
\indent Coefficientes aequationum per radicum potestates additas facile dantur.
\par\noindent\hfill[1796] Mai. 23.\quad Gött[ingae]

\medskip\begingroup\small
\indent Hierzu, wie auch zu der Nr. 28 vom 21. August 1796, vergl. man die oben S. 127, 128 abgedruckte Aufzeichnung aus \textsc{Leiste}.
\par\noindent\hfill\textsc{Klein}.
\par\endgroup
